\begin{abstract}
In a typical campus single sign-on deployment every service receives the same identifier and a fixed set of personal claims, no release is tied to a stated purpose, and students cannot see who received what. The survey by Werner, Westphall and Westphall shows that accountability and use limitation are among the least covered privacy characteristics in identity management. We present TrustTrail, a disclosure layer beside an unmodified OpenID Connect provider (Keycloak~26). Keycloak caps which of four disclosure profiles a service may obtain and issues per-service pairwise pseudonyms; a custom mapper removes the shared session identifier. TrustTrail is the only source of attribute values. It releases them only under a purpose declaration that the service has signed, which names the attributes, the purpose and a deletion date, and it returns a signed receipt. Releases and refused rule breaks are recorded in a hash-chained ledger whose Merkle roots are anchored in an Ethereum contract, and violations lower the service's tier automatically. On one machine, the layer added a median of 3.2\,ms per login (95\% CI 3.0--3.4\,ms). Conformance tests detected six tamper attacks and 45 scripted violations as designed, with no false positives in 90 honest releases. On the Sepolia test network, anchoring cost 254{,}408 gas per batch whatever its size. TrustTrail provides evidence of what services promised and received. It does not prove deletion, and its protection against the operator covers only history that has already been anchored.
\end{abstract}

\begin{IEEEkeywords}
Federated identity management, privacy, OpenID Connect, pairwise pseudonyms, purpose limitation, accountability, tamper-evident logging, Merkle tree, blockchain anchoring, trust management.
\end{IEEEkeywords}
